%-----------------------------------------------------------------------------------------------------------------------------------------------%
%	The MIT License (MIT)
%	Copyright (c) 2021 Jitin Nair
%-----------------------------------------------------------------------------------------------------------------------------------------------%

\documentclass[a4paper,12pt]{article}

%----------------------------------------------------------------------------------------
%	FONT (Enable XeLaTeX for Chinese)
%----------------------------------------------------------------------------------------
\usepackage{fontspec}
\usepackage{xeCJK}
\defaultfontfeatures{Ligatures=TeX}
\setmainfont{TeX Gyre Termes} % safe default on Overleaf
\setsansfont{TeX Gyre Heros}
\setmonofont{TeX Gyre Cursor}
% Noto CJK is what Overleaf ships; macOS has Songti SC instead. Probe at compile
% time so this one file builds unchanged in both places.
\IfFontExistsTF{Noto Serif CJK SC}{
  \setCJKmainfont{Noto Serif CJK SC}
  \setCJKmonofont{Noto Sans CJK SC}
}{
  \setCJKmainfont{Songti SC}
  \setCJKmonofont{Songti SC}
}

%----------------------------------------------------------------------------------------
%	PACKAGES
%----------------------------------------------------------------------------------------
\usepackage{url}
\usepackage{parskip}

\RequirePackage{color}
\RequirePackage{graphicx}
\usepackage[dvipsnames]{xcolor}
\usepackage[scale=0.9]{geometry}

\usepackage{tabularx}
\usepackage{enumitem}

\newcolumntype{C}{>{\centering\arraybackslash}X}

\usepackage{supertabular}
\usepackage{tabularx}
\newlength{\fullcollw}
\setlength{\fullcollw}{0.47\textwidth}

\usepackage{titlesec}
\usepackage{multicol}
\usepackage{multirow}

\titleformat{\section}{\Large\scshape\raggedright}{}{0em}{}[\titlerule]
\titlespacing{\section}{0pt}{10pt}{10pt}

% publications
\usepackage[style=authoryear,sorting=ynt,maxbibnames=10]{biblatex}

% hyperref
\usepackage[unicode, draft=false]{hyperref}
\definecolor{linkcolour}{rgb}{0,0.2,0.6}
\hypersetup{colorlinks,breaklinks,urlcolor=linkcolour,linkcolor=linkcolour}
\addbibresource{citations.bib}
\setlength\bibitemsep{0.7em}

% social icons
\usepackage{fontawesome5}

%----------------------------------------------------------------------------------------
% Job listing environments
%----------------------------------------------------------------------------------------
\newenvironment{jobshort}[2]
{
\begin{tabularx}{\linewidth}{@{}l X r@{}}
\textbf{#1} & \hfill &  #2 \\[3.75pt]
\end{tabularx}
}
{
}

\newenvironment{joblong}[2]
{
\begin{tabularx}{\linewidth}{@{}l X r@{}}
\textbf{#1} & \hfill &  #2 \\[3.75pt]
\end{tabularx}
\begin{minipage}[t]{\linewidth}
\begin{itemize}[nosep,after=\strut,leftmargin=1.2em,itemsep=3pt,label=--]
}
{
\end{itemize}
\end{minipage}
}

\usepackage{etaremune}
\usepackage{amssymb} % Required for the diamond symbol (\lozenge)

%----------------------------------------------------------------------------------------
%	PRIVATE DETAILS
%----------------------------------------------------------------------------------------
% The phone number is omitted by default, because the default build is the one
% CI publishes to the public site. `make private` defines \PrivateBuild as 1 on
% the command line to include it, for copies sent directly to people.
\providecommand{\PrivateBuild}{0}

%----------------------------------------------------------------------------------------
%	BEGIN DOCUMENT
%----------------------------------------------------------------------------------------
\begin{document}
\pagestyle{empty}

%----------------------------------------------------------------------------------------
%	TITLE
%----------------------------------------------------------------------------------------
\begin{tabularx}{\linewidth}{@{} C @{}}
{\Huge Xihan Yao \quad {\Large }} \\[6pt]

Austin, TX \ $|$ \
\ifnum\PrivateBuild=1
\href{tel:+13413339583}{\raisebox{-0.05\height}\faMobile\ +1 (341) 333 9583} \ $|$ \
\fi
\href{mailto:xyaoaf@utexas.edu}{\raisebox{-0.05\height}\faEnvelope\ xyaoaf@utexas.edu} \ $|$ \
\href{https://www.linkedin.com/in/xihan-yao-6381b3181/}{\raisebox{-0.05\height}\faLinkedin\ LinkedIn} \ $|$ \
\href{https://github.com/xyaoaf}{\raisebox{-0.05\height}\faGithub\ GitHub} \ $|$ \
\href{https://xyaoaf.github.io/}{\raisebox{-0.05\height}\faGlobe\ Website}

\end{tabularx}



%----------------------------------------------------------------------------------------
%	RESEARCH SUMMARY
%----------------------------------------------------------------------------------------
\section{Summary}
GeoData/Remote Sensing researcher building reproducible, large-scale geospatial analytics.
Experienced in land use/land cover modeling, spatial feature extraction, and multi-modal urban data (Street View, land cover, POI, mobility). Applies GeoAI and statistical modeling to study human–environment interactions, with strong Python-based pipelines and spatial diagnostics.


%----------------------------------------------------------------------------------------
%	EDUCATION
%----------------------------------------------------------------------------------------
\section{Education}
\begin{tabularx}{\linewidth}{@{}l X@{}}
2025 -- present &
\textbf{The University of Texas at Austin} --- PhD student, Geography and the Environment \\
& Supervisor: Professor Yuhao Kang \\
& Link: \href{https://sites.utexas.edu/gisense/}{GISense Lab}\\[3pt]

2021 -- 2023 &
\textbf{University of California, Berkeley} --- MLA, Environmental Planning \\[3pt]
& Thesis: Towards more effective urban vegetation monitoring and management: Studying of urban tree characteristics and
their relationship to the urban heat island effect with high-resolution remote sensing products, a case study in
Portland, Oregon, USA. \\
& Link: \href{https://doi.org/10.1016/j.landurbplan.2025.105420}{Publication} \\[3pt]

2017 -- 2021 &
\textbf{Hong Kong University of Science and Technology (HKUST)} --- BSc, Environmental Management and Technology \\[3pt]

2019 &
\textbf{University of Washington, Seattle} --- Exchange Student \\
\end{tabularx}

%----------------------------------------------------------------------------------------
%	EXPERIENCE
%----------------------------------------------------------------------------------------
\section{Experience}

\begin{joblong}{The University of Texas at Austin --- Research Assistant}{Austin, TX, 2025 -- present}
\item Leading a preregistered, national-scale research project modeling relationships between built environment characteristics and Big Five personality traits across U.S. urbanized areas.
\item Designing reproducible, Python-based workflows to process and harmonize TB-scale, multi-modal geospatial data, including Google Street View imagery, high-resolution land cover, POI data, and human mobility records.
\item Applying Earth foundation models (e.g., Google AlphaEarth) and statistical modeling to evaluate whether high-capacity embeddings provide additional explanatory power beyond traditional built-environment variables.
\item Conducting model diagnostics and spatial analyses (residual analysis, spatial autocorrelation tests) to identify geographic outliers, systematic failure modes, and limits of model generalizability.
\item Supervising four undergraduate research assistants and coordinating collaborative research across the lab and with external collaborators.
\end{joblong}


\begin{joblong}{EarthDefine, LLC --- GIS Analyst}{Redmond, WA, 2024 -- 2025}
\item Designed and implemented deep learning–based raster classification workflows for city-scale, sub-meter land cover mapping and change detection, emphasizing reproducible batch processing and quality control.
\item Led a team of GIS specialists to deliver operational geospatial data products across the U.S., including community-level canopy coverage metrics and multi-year land-use change datasets; supported external stakeholders such as the Arbor Day Foundation and municipal governments.
\end{joblong}

%----------------------------------------------------------------------------------------
% COMBINED UC BERKELEY EXPERIENCE
%----------------------------------------------------------------------------------------
\textbf{University of California, Berkeley} \hfill Berkeley, CA, 2022 -- 2023 \\
\vspace{-18pt}

\begin{itemize}[leftmargin=1.5em, label={}]
    % Role 1: Lecturer
    \item \textbf{— Lecturer, GIS} \hfill 2023
    \begin{itemize}[nosep, leftmargin=1.2em, label=--]
        \item Served as instructor of record for UC Berkeley’s largest GIS course (GEOG/LDARCH C188; 200 students), delivering lectures, designing assessments, and supervising a multi-TA instructional team.
        \item Redesigned and modernized the curriculum to incorporate contemporary GIS concepts and tools, working with ESRI to improve data workflows, reproducibility, and instructional data management.
    \end{itemize}
    
    \vspace{3pt}

    % Role 2: Graduate Student Instructor
    \item \textbf{— Graduate Student Instructor, Applied Remote Sensing \& GIS} \hfill 2022 -- 2023
    \begin{itemize}[nosep, leftmargin=1.2em, label=--]
        \item Taught applied remote sensing labs using Google Earth Engine (JavaScript), covering object- and pixel-based image analysis, machine-learning land use classification, and multi-source data integration (Landsat, MODIS, NAIP, LiDAR).
        \item Led hands-on GIS labs in ArcGIS Pro, updated legacy course materials, and produced instructional videos to support the transition from ArcGIS Desktop to ArcGIS Pro.
    \end{itemize}
\end{itemize}

% \begin{joblong}{Beijing Smart \& Green Transport Technology --- Data Research and Analysis Intern}{2020}
% \item Analyzed carbon abatement potential in China's transportation sectors under various green technology incentive policy scenarios; utilized a multi-parameter model to assess effectiveness under the supervision of Dr. Dongquan He.
% \end{joblong}
%----------------------------------------------------------------------------------------
%	AWARDS \& CERTIFICATES
%----------------------------------------------------------------------------------------
\section{Awards}

\begin{itemize}[leftmargin=1.5em, itemsep=4pt, label=--]

    \item \textbf{Energy AI Hackathon (6th Annual)}, UT Austin — Hildebrand Department of Petroleum and Geosystems Engineering (PGE). \\
    Awarded \textbf{Fourth Place} in a competitive machine learning challenge (46 teams). 
    Developed a "similarity–geospatial proximity gated model" achieving strong predictive accuracy and interpretability for subsurface production forecasting under sparse, heterogeneous conditions. \hfill 2026

    \item \textbf{Geoscience Hackathon}, UT Austin — Jackson School of Geosciences. \\
    Awarded \textbf{Second Place}. Project: ``Tracking Earth's Changes with AlphaEarth Foundations,'' applying foundation model embeddings for large-scale, flexible landscape system change detection. \hfill 2025

    \item \textbf{CyberTraining Summer School}, Texas A\&M University — Department of Geography. \\
    Awarded \textbf{First Place}. Project: ``Assessing Instance-based Disaster Impact Through Multimodal Geospatial Data''  \hfill 2025

    \item \textbf{IEEE GRSS GeoPitch Award}, IEEE Geoscience and Remote Sensing Society. \\
    Awarded \textbf{Third Place}. Elevator Pitch Presentation: ``Object-Based Urban Trees Characterization with Airborne LiDAR'' \hfill 2023

\end{itemize}


%----------------------------------------------------------------------------------------
%	PUBLICATIONS (BibLaTeX)
%----------------------------------------------------------------------------------------
\section{Selected Publications}

\begin{etaremune}[leftmargin=1.5em, itemsep=0.7em, start=2]
    % Publication 3 (Newest)
    \item \textbf{Yao, X.}, Kim, M., Dronova, I., McBride, J. R., Kondolf, G. M., \& Radke, J. D. (2025). Community-scale microclimate simulation using Airborne Laser Scanning and object-based urban tree classification. \textit{Landscape and Urban Planning, 263}, 105420. DOI: \href{https://doi.org/10.1016/j.landurbplan.2025.105420}{10.1016/j.landurbplan.2025.105420}

    % Publication 2
    \item Tong, J. C. K., \textbf{Yao, X. H.}, Lau, E. S. L., Cheung, W. K. S., Ho, K. K. S., Ng, V. Y. Y., \& Lau, A. P. S. (2023). Light pollution impact assessment in Hong Kong: Multi-dimensional measurement and spatial numerical modelling on integrated light sources in the neighbourhood level. \textit{Energy \& Environment, 35}(5), 2497--2516. DOI: \href{https://doi.org/10.1177/0958305X221146942}{10.1177/0958305X221146942}

\end{etaremune}


%----------------------------------------------------------------------------------------
%	SKILLS
%----------------------------------------------------------------------------------------
\section{Skills}
\begin{tabularx}{\linewidth}{@{}l X@{}}
\textbf{Programming \& Data Science} & 
Python (TensorFlow, PyTorch, Scikit-learn, GeoPandas, Rasterio, PySAL), R, JavaScript (Google Earth Engine) \\ [3pt]
\textbf{GIS \& Remote Sensing} & 
ESRI products (ArcGIS Pro, ArcGIS Online), QGIS, Google Earth Engine, Google AlphaEarth Foundations, LiDAR \\ [3pt]
\textbf{Computing \& Systems} & 
Linux Shell, High-Performance Computing (HPC), GPU server management, batch processing, Git/GitHub \\
\end{tabularx}

\vfill
\center{\footnotesize Last updated: \today}

\end{document}
